\begin{center}
{\LARGE\bfseries Population gradient flow at fixed depth}\\[5pt]
{\large Why a positive physical-time interval survives the width limit}\\[7pt]
{\small Main exposition with a separate technical appendix \quad\textbullet\quad 6 September 2026}
\end{center}

\section*{1. The network and the question}

The conclusion is a population law for training through a \emph{fixed positive
physical time} $T_*$. Width tends to infinity and the gradient-descent step
tends to zero together, with no restriction linking their rates. The interval
can depend on depth and on fixed bounds describing the model; it does not
shrink with width or step size. This is the sense of ``semi-global'' intended
here. In the usual terminology it is a local-in-time theorem, not an assertion
about every finite horizon.

There are $L\ge2$ hidden layers of width $n$, and $m$ fixed inputs
$x_a\in\mathbb R^d$ with labels $y_a$. The indices $a,b$ denote examples,
$i,j$ neurons, $\ell$ layers, and $k$ training steps. Both $L$ and the dataset
are fixed as $n$ grows. Each layer may have its own coordinatewise activation
$\phi^{(\ell)}$. The finite forward pass is
\[
\begin{aligned}
z_a^{(1)}&=W^{(1)}x_a/\sqrt d,
&h_a^{(1)}&=\phi^{(1)}(z_a^{(1)}),\\
z_a^{(\ell)}&=W^{(\ell)}h_a^{(\ell-1)},
&h_a^{(\ell)}&=\phi^{(\ell)}(z_a^{(\ell)}),\quad 2\le\ell\le L,\\
f_{n,a}&=(W^{(L+1)})^\top h_a^{(L)}/n.
\end{aligned}
\]
$W^{(1)}$ is $n$ by $d$, each middle $W^{(\ell)}$ is $n$ by $n$, and
$W^{(L+1)}$ is an $n$-vector. The last object is always the \emph{stored
rescaled readout}. In particular, when $L=2$, the prediction is
$(W^{(3)})^\top h_a^{(2)}/n$.

The backpropagated fields contain the derivative of the prediction, but
\emph{not} the loss residual:
\[
\begin{aligned}
\delta_a^{(L)}&=W^{(L+1)}\phi^{(L)\prime}(z_a^{(L)}),\\
\delta_a^{(\ell)}&=\phi^{(\ell)\prime}(z_a^{(\ell)})
       (W^{(\ell+1)})^\top\delta_a^{(\ell+1)},\qquad \ell<L.
\end{aligned}
\]
Products with $\phi^{(\ell)\prime}$ are coordinatewise. Equivalently,
$\delta_a^{(\ell)}=n\,\partial f_{n,a}/\partial z_a^{(\ell)}$.

We use the weighted squared loss
\[
\mathcal L_n=\sum_{a=1}^m\omega_a(f_{n,a}-y_a)^2,
\qquad \omega_a>0,\qquad\sum_{a=1}^m\omega_a=1,
\]
and stored-weight learning rates
\[
\eta_n\bigl(n\kappa_1,\kappa_2,\ldots,\kappa_L,n\kappa_{L+1}\bigr),
\qquad t=k\eta_n.
\]
The fixed $\kappa_\ell\ge0$ set the relative training speeds of the layers.
The common $\eta_n$ resolves physical time more finely; it is not present
in the final differential equation.

\newpage
\section*{2. The configuration covered by the theorem}

\textbf{Activations.} Each $\phi^{(\ell)}$ is continuously differentiable,
its derivative is bounded, and its derivative is globally Lipschitz.
The constants bounding these quantities and $|\phi^{(\ell)}(0)|$ are fixed.
Activation values themselves may be unbounded: bounded derivatives imply
at most linear growth. Arctan, tanh, sigmoid, softplus, GELU and SiLU are
examples. No oddness, monotonicity, analyticity or positive lower bound on
the derivative is required. ReLU itself is outside this result.

\textbf{Initialization.} First-weight entries are independent, identically
distributed, centered subGaussian variables. Each middle matrix has
independent entries
\[
W_{0,ji}^{(\ell)}\sim
\mathcal N(0,\sigma_\ell^2/n),\qquad 2\le\ell\le L.
\]
Readout entries are independent copies of a fixed subGaussian variable;
their mean need not vanish. All initialization groups are independent.
Here subGaussian means that absolute moments grow no faster than a fixed
constant times $\sqrt p$ in $L^p$, or equivalently that tails have a Gaussian
upper bound. Zero readout, bounded readout and order-one Gaussian readout
are all allowed. Gaussian first weights are a special case. Zero variance
or a zero learning multiplier is allowed for existence.

\textbf{Data.} The normalized input lengths $\|x_a\|_2/\sqrt d$ and labels
are bounded. Inputs may coincide; their Gram matrix may be singular. Its
entries will be written explicitly as $x_a^\top x_b/d$. No inverse Gram
matrix or angular separation is used in the existence proof.

\medskip
\textbf{Theorem (fixed positive-time joint limit).}
For this configuration there is $T_*>0$ such that, for \emph{every}
$\eta_n\to0$, gradient descent has the population limit described below,
throughout $[0,T_*]$, in the observable senses stated here. Parameters are
interpolated linearly, and forward quantities are recomputed from them.
The population integral equations have a unique strong solution on their
realized field and operator spaces. Moreover:
\begin{itemize}
\item predictions, loss and all tangent-kernel entries converge in probability,
uniformly in time;
\item in each hidden layer separately, the empirical joint laws of the
examples' preactivation and activation paths converge, including second
moments measured with the uniform path norm;
\item the time integrals of squared hidden-coordinate speeds converge.
\end{itemize}
``Strong'' here means that fields satisfy the equations in mean square and
the middle operators in operator norm. It does not mean coordinatewise
uniform control over all neurons.

$T_*$ depends only on fixed depth, activation bounds, initialization scales,
input/label bounds and layer learning multipliers. Under normalized loss
weights, it can be chosen independently of $m,d$ and the input angles when
these bounds remain fixed. The convergence assertion still fixes $m,d$ and
the dataset: this is not a theorem for datasets growing with width. For the
unweighted \emph{summed} loss, physical time is multiplied by $m$ relative
to the uniformly weighted loss. No interval uniform over increasing depth
is claimed.

\newpage
\section*{3. What survives in the population description}

There is one neuron population for each hidden layer. Uppercase
$Z_a^{(\ell)},H_a^{(\ell)}$ describe a coordinate in that population, jointly
across examples and times. They retain a distribution: taking width to
infinity makes empirical averages deterministic, not every neuron identical.
An expectation below averages within the population containing its factors.

Each $W^{(\ell)}$, $2\le\ell\le L$, becomes a bounded linear map from
mean-square fields in layer $\ell-1$ to those in layer $\ell$. The reverse
map is its adjoint $(W^{(\ell)})^*$, inherited from the \emph{same} finite
matrix. It is not an independently sampled backward operator. Thus
\[
\begin{aligned}
H_a^{(1)}&=\phi^{(1)}(Z_a^{(1)}),\\
Z_a^{(\ell)}&=W^{(\ell)}H_a^{(\ell-1)},\qquad
H_a^{(\ell)}=\phi^{(\ell)}(Z_a^{(\ell)}),\quad 2\le\ell\le L,\\
f_a&=\mathbb E[W^{(L+1)}H_a^{(L)}],\\
\delta_a^{(L)}&=W^{(L+1)}\phi^{(L)\prime}(Z_a^{(L)}),\\
\delta_a^{(\ell)}&=\phi^{(\ell)\prime}(Z_a^{(\ell)})
             (W^{(\ell+1)})^*\delta_a^{(\ell+1)},\quad\ell<L.
\end{aligned}
\]
To see the middle-layer evolution, begin with its exact finite update:
\[
W_{k+1}^{(\ell)}-W_k^{(\ell)}
=-\frac{2\eta_n\kappa_\ell}{n}
\sum_b\omega_b(f_{n,b,k}-y_b)
\delta_{b,k}^{(\ell)}(h_{b,k}^{(\ell-1)})^\top.
\]
Applied to a vector, this is a sum of backward fields multiplied by
feature overlaps. A finite overlap
$(h_a^{(\ell-1)})^\top h_b^{(\ell-1)}/n$ becomes the population average
$\mathbb E[H_a^{(\ell-1)}H_b^{(\ell-1)}]$.
For any fixed mean-square test field $V$ in layer $\ell-1$, the complete
evolution is therefore
\[
\begin{aligned}
\dot Z_a^{(1)}
 &=-2\kappa_1\sum_b\omega_b\frac{x_a^\top x_b}{d}
          (f_b-y_b)\delta_b^{(1)},\\
\dot W^{(\ell)}V
 &=-2\kappa_\ell\sum_b\omega_b(f_b-y_b)
          \delta_b^{(\ell)}\mathbb E[H_b^{(\ell-1)}V],\quad 2\le\ell\le L,\\
\dot W^{(L+1)}
 &=-2\kappa_{L+1}\sum_b\omega_b(f_b-y_b)H_b^{(L)}.
\end{aligned}
\]
The derivative in the middle line acts on $W^{(\ell)}$, with $V$ fixed.
The initial operators are defined by joint limits of all the needed forward
and transpose computations. Their construction, and passage of the finite
adjunction identity to the limit, are given in the appendix.

These equations use only the current state. Earlier training is stored in
the current operators and fields; a separate history is not required to
evolve them. This is an autonomous population law with infinitely many
scalar degrees of freedom, not a finite-dimensional closure.

\newpage
\section*{4. The mechanism of feature learning and the proof's task}

The population equations keep the nonlinear forward and backward passes.
Their content is especially visible after differentiating a hidden
preactivation, for $2\le\ell\le L$:
\[
\dot Z_a^{(\ell)}
=\dot W^{(\ell)}H_a^{(\ell-1)}
 +W^{(\ell)}\bigl[
 \phi^{(\ell-1)\prime}(Z_a^{(\ell-1)})\dot Z_a^{(\ell-1)}\bigr].
\]
The first term changes the receiving layer's weights. The second transports
the movement of the preceding layer through its current weights. Thus
training can modify both what a layer receives and how it processes that
input. Neither term has a vanishing width factor under the stored scaling.

Let $K_{ab}(t)$ be the current tangent-kernel entry, including the fixed
$\kappa_\ell$. It is the sum of parameter-gradient inner products from
the layers; its explicit contractions are in the appendix. The output law is
\[
\dot f_a=-2\sum_b\omega_bK_{ab}(t)(f_b-y_b),
\]
and
\[
\frac{d\mathcal L}{dt}
=-4\sum_{a,b}\omega_a\omega_b(f_a-y_a)K_{ab}(t)(f_b-y_b)
\le0.
\]
The sign follows because $K(t)$ is a Gram matrix with nonnegative layer
weights. The kernel is recomputed from the evolving features and backward
fields. The theorem does not freeze it at initialization or replace the
activations by affine functions.

Existence does not force strict movement in every layer: zero training
multipliers, zero residuals and degenerate configurations are permitted.
With zero initial readout, all hidden backpropagated fields initially vanish,
although the readout can start moving immediately. Subsequent readout
movement creates backward signals that drive hidden learning. The separate
strict-activity result for two hidden layers quantifies this onset, under its
additional nondegeneracy assumptions. Such a result for every depth is not
part of the present theorem.

The convergence proof has to justify the full feedback loop, not just write
its formal differential equation. It proceeds through four concrete tasks:
\begin{enumerate}
\item Bound typical field sizes and matrix amplification for a common
positive time, at every width and every sufficiently small step.
\item Account for reuse of the initial Gaussian matrices, then control the
tails of the backward signals uniformly as the time mesh is refined.
\item Use those tails to compare different time discretizations, constructing
and identifying a unique continuous population flow.
\item Compare actual finite-width training to a fixed, coarser reference
computation, and then remove that reference mesh.
\end{enumerate}
The middle two tasks contain the main analysis. A width limit for a fixed
number of operations does not on its own control the growing number
$T_*/\eta_n$ of gradient-descent steps.

\newpage
\section*{5. Why a fixed time is possible, and why size bounds are not enough}

The correct notion of a typical finite field size is
$\|h_a^{(\ell)}\|_2/\sqrt n$ or $\|\delta_a^{(\ell)}\|_2/\sqrt n$.
In the population it is the square root of the corresponding second moment.
A middle matrix is measured by its ordinary operator norm, which controls
its amplification of these field sizes.

These norms match the training update exactly:
\[
\left\|\frac{\delta_a^{(\ell)}(h_a^{(\ell-1)})^\top}{n}
\right\|_{\operatorname{op}}
=\frac{\|\delta_a^{(\ell)}\|_2}{\sqrt n}
 \frac{\|h_a^{(\ell-1)}\|_2}{\sqrt n}.
\]
The factor $1/n$ balances two vectors of Euclidean size $\sqrt n$.
As long as field sizes, readout size and matrix norms stay in a fixed bounded
range, all parameter velocities are bounded by constants independent of
width. Forward propagation uses the activation's linear growth; backward
propagation uses its bounded derivative and finitely many matrix norms.
Weighted training sums contribute total weight one.

Initialization lies in such a range with probability tending to one. Choose
a larger range with a fixed margin. Before a hypothetical first exit, the
total motion is at most the uniform velocity bound times elapsed time.
Choosing time smaller than margin divided by that bound prevents the exit.
The same argument applies to discrete Euler updates and population integral
equations. This supplies the first positive interval. Its constants may grow
with depth, but contain neither $n$ nor the number of time steps.

Why does this not already prove the theorem? The backward pass contains
multiplication by an unbounded random field. At an internal layer this is
\[
\phi^{(\ell)\prime}(Z_a^{(\ell)})
       (W^{(\ell+1)})^*\delta_a^{(\ell+1)}.
\]
A change in $Z_a^{(\ell)}$ multiplies the existing backward signal. A
Lipschitz derivative bounds that change by a constant times
$|Z_a^{(\ell)}-\widetilde Z_a^{(\ell)}|$, where the tilde denotes a comparison
state. But a second-moment bound on each
factor does not bound the second moment of their product. The ordinary
finite-dimensional Lipschitz proof would instead use the largest coordinate,
which can grow with width. That would spoil a uniform time estimate.

There is also a dependence issue. For example,
$(W^{(\ell)})^\top W^{(\ell)}h_a^{(\ell-1)}$ uses the same matrix twice.
Conditioning on earlier forward and backward calls changes the law of later
calls. Treating every call as independent Gaussian noise discards the very
correlations generated by feature learning. The proof must preserve these
correlations while controlling the rare, large backward signals that make
the comparison difficult.

\newpage
\section*{6. Matrix reuse, small feedback pulses, and uniform tails}

For a fixed Euler mesh $\Delta$, each trained middle matrix equals its initial
Gaussian matrix plus a finite sum of rank-one updates. The learned terms
therefore involve only previously computed features, backward fields and
their scalar overlaps. The remaining task is to understand repeated calls
to the same initial Gaussian matrix and its transpose.

The fixed-computation Gaussian response theorem supplies exactly this
description \cite{TP3}: a new call has a Gaussian part together with
correction terms recording its dependence on previous calls in the opposite
direction. The correction coefficients are expectations of sensitivities
to the Gaussian inputs in those earlier calls. Their covariance laws,
residuals and scalar coefficients are held fixed in taking these
sensitivities. The Gaussian parts at different times are correlated; neither
time independence nor finite-width neuron independence after reuse is used.

The decisive estimate is more precise than a bound on the total response.
A backward Gaussian input belonging to example $b$ at an earlier step can
affect a later forward feature only through training. That first change
carries the update weight $\Delta\omega_b$. Propagating the change through
later steps preserves this factor:
\[
\bigl|\text{expected forward sensitivity to that backward input}\bigr|
\le C\Delta\omega_b.
\]
Here $C,c>0$ are fixed constants determined by the model bounds. The appendix
proves these sensitivity estimates, including the reverse-direction bounds.

This small pulse is what turns sums over many training steps into integrals
over a fixed amount of time. The random amplification involves expressions
of the form
\[
\Delta\sum_{s<k}\sum_b\omega_b
\left|(W_s^{(\ell+1)})^*\delta_{b,s}^{(\ell+1)}\right|,
\]
where $s$ indexes earlier steps. It does \emph{not} require the maximum of
these random fields over a refining time grid. Jensen's inequality applied
to the weights $\Delta\omega_b/(k\Delta)$ controls the exponential of this
time average using marginal Gaussian-tail bounds, even with complete
dependence across times and inputs. The weights sum to one: the exponential
of their average is at most the corresponding average of exponentials.

The size bound controls every Gaussian innovation's variance. Forward fields
are a Gaussian or root contribution plus time-weighted backward fields;
backward fields are Gaussian contributions plus response-weighted forward
fields. Their tail estimates must therefore be closed together.

To do this without circular reasoning, choose forward-response bounds from the first
layer upward, using their zero-time prefactors. Choose backward-response
bounds from the readout downward in the same way. At zero time the opposite
direction enters only through terms multiplied by elapsed time. Once all
these finitely many bounds are fixed, reduce $T_*$ so those terms are small.
The forward-then-backward construction at each Euler step preserves the
bounds; it never assumes a future estimate to prove a current one.

The resulting internal-layer estimate is
\[
\sup_{k\Delta\le T_*}\mathbb E\exp\!\left(
c\left|(W_k^{(\ell+1)})^*\delta_{a,k}^{(\ell+1)}\right|^2\right)
\le C,
\]
uniformly in $\Delta$; the readout obeys the analogous bound. This controls
each time marginal, not the maximum of a random path. It is the additional
estimate that makes the continuous-time argument possible.

\newpage
\section*{7. Localized stability constructs the flow and proves uniqueness}

Compare two states on the same population spaces. Measure their difference
by the first-preactivation and readout mean-square distances, together with
the middle-operator norm distances. Forward propagation is Lipschitz in
this distance on the bounded range from page 5.

For the backward pass, split the difference into two effects: a change in
the incoming backward field multiplied by a bounded activation derivative,
and a change in the derivative multiplied by the reference backward field.
For the second effect, choose a cutoff $R>0$.

On coordinates where the absolute reference backward signal is at most $R$,
the contribution is bounded by a constant times
\[
R\,|Z_a^{(\ell)}-\widetilde Z_a^{(\ell)}|.
\]
On the remaining coordinates, the preceding Gaussian-tail estimate makes
its mean-square norm at most $Ce^{-cR^2}$. This estimate applies to the
readout at the final layer as well, so a bounded readout is not needed.

The way this comparison passes through depth matters. At each preceding
layer, the error already present is multiplied only by a bounded derivative
and a bounded operator. The layer adds a new cutoff contribution. It does
not multiply the entire incoming error by $R$. Thus the final stability
constant grows like $C(1+R)$, with depth absorbed into $C$, rather than
like $R^L$.

Apply this comparison to two population Euler paths with meshes $\Delta$
and $\Delta'$. Their grid states differ from their interpolated states by
at most a constant times their mesh lengths. Integration and Gronwall's
inequality give
\[
\begin{split}
\text{largest state discrepancy on }[0,T_*]
\ \le\ & C e^{C(1+R)T_*}\\[-3pt]
&\times\left[(1+R)(\Delta+\Delta')+e^{-cR^2}\right].
\end{split}
\]
First hold $R$ fixed and let both meshes tend to zero. Then let $R$ grow.
The remaining term vanishes because
\[
e^{C(1+R)T_*-cR^2}\longrightarrow0.
\]
This balance is the heart of the proof: the amplification cost is exponential
in $R$, while the discarded tail decays exponentially in $R^2$.

The Euler paths are consequently Cauchy in the complete field and operator
spaces. Their limit satisfies the population integral equations; continuity
of bounded-derivative products is justified by a reference-field cutoff.
The tail estimates pass to this solution. If another strong solution starts
from the same state, use the constructed solution as reference in the same
comparison. The other solution needs only the bounded-range estimate.
Letting $R$ grow forces their distance to zero. This proves uniqueness,
without assuming Gaussian tails for every possible competing solution.

\newpage
\section*{8. Why every vanishing gradient-descent step is allowed}

We still have to connect the population construction to actual training,
which takes approximately $T_*/\eta_n$ steps. Applying the fixed-computation
theorem directly to this growing program would not be justified.

Instead fix a coarse mesh $\Delta$ independently of $n$ and $\eta_n$.
Compute its population Euler path. In a finite-width reference computation,
use the \emph{same sampled initial arrays} as actual gradient descent,
but supply the scalar residuals and feature/backward overlaps from that
population Euler path. After expanding the matrices into their initial
parts and trained rank-one terms, this is a fixed finite computation.
The Gaussian response theorem therefore proves empirical convergence of
all its required fields and scalar overlaps.

There is a necessary consistency step. Reassemble actual finite matrices
from the reference rank-one updates and recompute their forward and backward
passes. These need not equal the prescribed reference calls exactly:
recomputed overlaps are empirical averages, while the prescribed calls use
their population expectations. At the fixed mesh there are only finitely
many such differences. They vanish as width grows. The backward comparison
uses reference tails again to handle the unbounded readout. This produces
a finite parameter path that is an approximate solution of the \emph{actual}
training equations, with a vanishing error at fixed $\Delta$.

Now compare actual fine-step gradient descent to this coarse reference.
Both have uniform field-size and operator bounds. Only the reference needs
the tail estimate. The same localized stability argument gives
\[
\begin{split}
\text{largest finite-state discrepancy on }[0,T_*]
\ \le\ & C e^{C(1+R)T_*}\\[-3pt]
&\times\left[(1+R)(\eta_n+\Delta)+e^{-cR^2}
                 +o_{\mathbb P}(1)\right].
\end{split}
\]
Here $o_{\mathbb P}(1)$ denotes a quantity tending to zero in probability
as $n\to\infty$, with $R,\Delta$ fixed. The state comparison is between
two finite systems on the same sampled arrays; no operator-norm comparison
between a finite matrix and a population operator is asserted.

The order of limits is deliberate. First let $n\to\infty$ with the desired arbitrary sequence $\eta_n\to0$,
while $R,\Delta$ stay fixed. Then refine $\Delta\to0$ using the population
Euler convergence from page 7. Finally send $R\to\infty$. Each contribution
to the displayed error vanishes in that order.

The coarse computation is used only in the proof. It does not replace the
feedback of actual training, nor become an external forcing in the population
law. Crucially, its number of operations stays fixed during the width limit.
There is therefore no requirement such as $\eta_n\sqrt n\to0$.

\newpage
\section*{9. Completing the result and locating the technical details}

Once state comparisons are available, predictions follow from the forward
Lipschitz bounds. Kernel entries are products of same-population second
moments; convergence of forward and backward fields controls those moments
by Cauchy--Schwarz. There is no need for a central limit theorem or an
independence claim about trained neurons.

To recover whole hidden paths, first record a finite grid of times. The
fixed-computation limit identifies their joint empirical laws. Between
neighboring grid points, the squared variation of each coordinate is bounded
by the cell length times its integrated squared speed. Averaging this
inequality over neurons controls the error of reconstructing the paths from
the grid. Refining the grid therefore gives convergence with second moments
in the uniform path norm. The appendix carries out the additional cutoff
argument for convergence of integrated squared speeds.

The proof is first written for twice differentiable activations with bounded
first and second derivatives. Smooth approximations to a continuously
differentiable activation with Lipschitz derivative preserve these same
bounds. The time interval and stability constants are consequently uniform
over the smoothing scale. Passing that scale to zero gives the stated
activation class. This passage does not cover a jump in the derivative,
which is why ReLU is excluded.

More general separable losses are permitted. Write them as
\[
\mathcal L=\sum_a\omega_a\ell_a(f_a),
\]
where each loss derivative is locally Lipschitz, with common finite bounds
and Lipschitz constants on bounded prediction intervals. Replace every
$2(f_b-y_b)$ in the updates by $\ell_b'(f_b)$. The size argument uses its
bound; the feedback comparison uses its Lipschitz constant. No convexity is
needed for the fixed positive-time result.

Vanishing initialization changes are also allowed: first preactivations and
readout must change by a vanishing mean-square amount, and middle matrices
by a vanishing operator norm. In particular, zero population readout includes
finite Gaussian readouts whose stored coordinate scale is $n^{-\beta}$ for
any fixed $\beta>0$. This does not imply universality for order-one changes
to the middle-matrix distribution.

The appendix supplies the exact finite updates and kernel, the common
operator construction, the complete Gaussian response recursions, the
weighted sensitivity induction and choice of the depth-dependent constants,
the comparison estimates, and the passages for paths, losses and smoothing.
The main mechanism is preserved above: \emph{size bounds give room to evolve;
the update-weighted response controls tails despite matrix reuse; the tails
make the nonlinear feedback stable; a fixed coarse computation connects
that stable flow to every vanishing-step finite training sequence.}

The stored scaling, fixed-step width limits and operator viewpoint have prior
literature \cite{TP3,TP4,TP4b,TP5}. The additional argument presented here
is the control uniform in time mesh and its use in the joint limit. This
exposition makes no publication-priority claim. It establishes a fixed
positive interval at each fixed depth, not all-time control, strict feature
activity at every depth, or reconstruction from initialization jets by Borel
or Pad\'e summation.
